\section{Fast LRALM}
\label{sec:acceleration}

At iteration \(k\), we take primal and dual gradient steps on
\[
    \min_{\bx\in\RR^n}\max_{\blmbd\in\RR^m}
    \mathcal P_s(\bx,\blmbd;\bz^k,\zero)
    =\min_{\bx\in\RR^n}\max_{\blmbd\in\RR^m}
    \left\{
    \calD_\rho(\bx,\blmbd)
    +\frac{1}{2\gamma}\|\bx-\bz^k\|^2
    -\frac{1}{2s}\|\blmbd\|^2
    \right\}.
\]
The dual step uses Nesterov's accelerated gradient
method~\cite{nesterov2018lectures}. For maximizing an \(L\)-smooth and
\(\mu\)-strongly concave function \(\phi\), this method reads
\[
    \bar{\bu}^{k}=\bu^k+
    \frac{1-\sqrt{\mu/L}}{1+\sqrt{\mu/L}}(\bu^k-\bu^{k-1}),
    \qquad
    \bu^{k+1}=\bar{\bu}^{k}+\frac1L\nabla\phi(\bar{\bu}^{k}).
\]
The function \(\blmbd\mapsto\mathcal P_s(\bx,\blmbd;\bz,\zero)\) is
\(s^{-1}\)-strongly concave. Taking \(\mu=s^{-1}\) and replacing \(L^{-1}\) by
the dual stepsize \(\tau\) gives the effective condition number and
extrapolation coefficient
\[
    \kappa:=\frac{s}{\tau},
    \qquad
    \omega:=\frac{\sqrt{\kappa}-1}{\sqrt{\kappa}+1}.
\]

We refer to the resulting method as Fast LRALM (Algorithm~\ref{alg:b2-prime});
it lies outside \(\calC_\theta\).
% because it incorporates Nesterov
% extrapolation in the dual variable and evaluates both partial gradients of
% \(\mathcal P_s\) at the extrapolated multiplier \(\bar{\blmbd}^k\), whereas
% every method in \(\calC_\theta\) evaluates \(\nabla\mathcal P_s\) at the
% current multiplier \(\blmbd^k\) without extrapolation. Extrapolation is
% applied only to the multiplier sequence, while the primal variable is updated
% by a plain gradient step.
\begin{algorithm}[t]
    \caption{Fast LRALM}
    \label{alg:b2-prime}
    \small
    \begin{algorithmic}[1]
        \STATE \textbf{Input:} Parameters \(\rho,s,\gamma,\eta,\tau>0\),
        \(\alpha\in(0,1]\); \(\kappa:=s/\tau\) and
        \(\omega:=\frac{\sqrt{\kappa}-1}{\sqrt{\kappa}+1}\); iteration limit \(K\).
        \STATE \textbf{Initialize:} \(\bx^0=\bz^0=\zero\),
        \(\blmbd^{-1}=\blmbd^0=\zero\).
             \FOR{\(k=0,1,\ldots,K-1\)}
            \STATE \(\bar{\blmbd}^k=\blmbd^k+\omega(\blmbd^k-\blmbd^{k-1})\).
            \STATE \(\bx^{k+1}=\bx^k-\eta\Bigl(\nabla_{\bx}\calL_\rho\bigl(\bx^k,\by(\bx^k,\bar{\blmbd}^k),\bar{\blmbd}^k\bigr)+\tfrac{1}{\gamma}(\bx^k-\bz^k)\Bigr)\).
            \STATE \(\blmbd^{k+1}=\bar{\blmbd}^k+\tau\Bigl(\calA(\bx^{k+1})-\by(\bx^{k+1},\bar{\blmbd}^k)-\tfrac{1}{s}\bar{\blmbd}^k\Bigr)\).
            \STATE \(\bz^{k+1}=\bz^k+\alpha(\bx^{k+1}-\bz^k)\).
            \STATE \textbf{Output:}
            \(\bigl(\bx^{k+1},\,\by(\bx^{k+1},\bar{\blmbd}^k),\,\bq(\bx^{k+1},\bar{\blmbd}^k)\bigr)\).
        \ENDFOR
    \end{algorithmic}
\end{algorithm}

\subsection{Iteration Complexity}

Suppose that Assumptions~\ref{ass:smoothness} and~\ref{ass:h} hold, and let
\(\gamma>0\) satisfy \(\gamma L_{\calD,\bx}<1/2\). By
Proposition~\ref{prop:curvature-smoothness}, for every \(s>0\) and
\(\bz\in\RR^n\), \(\mathcal P_s(\cdot,\blmbd;\bz,\zero)\) is
\(\mu_\gamma\)-strongly convex and \(L_x\)-smooth, while
\(\mathcal P_s(\bx,\cdot;\bz,\zero)\) is \(s^{-1}\)-strongly concave and
\(L_\lambda\)-smooth. Consequently, the
value functions and associated optimizers
\begin{align}
    d_s(\blmbd;\bz)
    &:= \min_{\bx\in\RR^n}\mathcal P_s(\bx,\blmbd;\bz,\zero),
    &
    \bx_s(\blmbd;\bz)
    &:= \argmin_{\bx\in\RR^n}\mathcal P_s(\bx,\blmbd;\bz,\zero),
    \label{eq:dual-value}\\
    p_s(\bz)
    &:= \max_{\blmbd\in\RR^m} d_s(\blmbd;\bz),
    &
    \blmbd_s^\star(\bz)
    &:= \argmax_{\blmbd\in\RR^m} d_s(\blmbd;\bz)
    \label{eq:saddle-value}
\end{align}
are uniquely defined, and strong duality gives
\begin{equation}
    p_s(\bz)
    = \min_{\bx\in\RR^n}\max_{\blmbd\in\RR^m}
      \mathcal P_s(\bx,\blmbd;\bz,\zero).
    \label{eq:minimax-equality}
\end{equation}
Moreover, \(\mathcal P_s(\cdot,\cdot;\bz,\zero)\) admits the unique saddle
point \(\bigl(\bx_s^\star(\bz),\blmbd_s^\star(\bz)\bigr)\), where
\(\bx_s^\star(\bz):=\bx_s\bigl(\blmbd_s^\star(\bz);\bz\bigr)\).

Let \(\{(\bx^k,\blmbd^k,\bz^k)\}_{k\ge0}\) be generated by
Algorithm~\ref{alg:b2-prime}, and write its output as
\((\bx^{k+1},\by^{k+1},\bq^k)\), where
\[
    \by^{k+1}:=\by(\bx^{k+1},\bar{\blmbd}^k),
    \qquad \bq^k:=\bq(\bx^{k+1},\bar{\blmbd}^k).
\]
The following lemma bounds the KKT residuals of the outputs.
\begin{lemma}
\label{lem:b2-prime-residual-decomposition}
Suppose that Assumptions~\ref{ass:smoothness} and~\ref{ass:h} hold. Then,
for every \(k\ge0\), \(\bq^k\in\partial h(\by^{k+1})\) and
\[
\begin{aligned}
    \bigl\|\nabla f(\bx^{k+1})+J_{\calA}(\bx^{k+1})^\top\bq^k\bigr\|
    &\le\left(\left|\frac1\eta-\frac1\gamma\right|+L_{\calD,\bx}\right)\|\bx^{k+1}-\bx^k\|
    +\frac1\gamma\|\bx^{k+1}-\bz^k\|,\\
    \|\calA(\bx^{k+1})-\by^{k+1}\|
    &\le\frac1\tau\|\blmbd^{k+1}-\bar{\blmbd}^k\|
    +\frac1s\|\bar{\blmbd}^k\|.
\end{aligned}
\]
\end{lemma}

\begin{lemma}
\label{lem:b2-prime-multiplier-bounds}
Suppose that Assumption~\ref{ass:h} holds and
\(\red{\tau L_\lambda\le\frac12}\). Then, for every \(k\ge0\), it holds that
\[
    \max\left\{\|\blmbd^k\|,\|\bar{\blmbd}^k\|\right\}\le3C_\lambda L_h,
\] 
where the constant \(C_\lambda:=\max\left\{4,\sqrt{\frac{8\rho}{\tau}}\right\}\) is independent of \(s\).
\end{lemma}

For the descent analysis, we impose the following conditions, with parameters fixed
throughout each run:
\begin{equation}
\begin{gathered}
    \gamma L_{\calD,\bx}\le\frac{1}{10},\qquad \eta=\gamma,\qquad
    0<\tau\le\frac{\rho}{4},\qquad 14\gamma\tau G_{\calA}^2\le1,\\
    s\ge\max\{\rho,4\tau\},\qquad
    0<\alpha\le\frac{1}{6\sqrt{\kappa}}.
\end{gathered}
\label{eq:b2-prime-stepsizes}
\end{equation}

{\color{red}
The analysis rests on a Lyapunov function. For \(\bz\in\RR^n\) and \(k\ge0\),
we measure the inner subproblem error at the center \(\bz\) by
\[
    \mathcal E_k(\bz)
    :=p_s(\bz)-d_s(\blmbd^k;\bz)
    +\frac{1}{2s}\|\bw^k-\blmbd_s^\star(\bz)\|^2
    +\frac{\mu_\gamma}{4}\|\bx^k-\bx_s(\bar{\blmbd}^{k-1};\bz)\|^2
    +\frac{G_{\calA}^2}{4\mu_\gamma}\|\blmbd^k-\bar{\blmbd}^{k-1}\|^2,
\]
where \(\bw^k:=\blmbd^{k-1}+\sqrt{\kappa}(\blmbd^k-\blmbd^{k-1})\) and
\(\bar{\blmbd}^{-1}:=\blmbd^0\), and define
\[
    \mathcal W_k:=p_s(\bz^k)+\frac12\mathcal E_k(\bz^k).
\]
The sequence \(\{\mathcal W_k\}_{k\ge0}\) controls the three step increments
in Lemma~\ref{lem:b2-prime-residual-decomposition} through the following
one-step descent and initial gap bounds.

\begin{lemma}
\label{lem:b2-prime-potential-descent}
Suppose that Assumptions~\ref{ass:smoothness} and~\ref{ass:h} and
\eqref{eq:b2-prime-stepsizes} hold. Then, for every \(k\ge0\),
\[
    \mathcal W_{k+1}-\mathcal W_k
    \le -\frac{1}{128\gamma}\|\bx^{k+1}-\bx^k\|^2
        -\frac{1}{16\tau}\|\blmbd^{k+1}-\bar{\blmbd}^k\|^2
        -\frac{\alpha}{6\gamma}\|\bx^{k+1}-\bz^k\|^2.
\]
\end{lemma}

\begin{lemma}
\label{lem:b2-prime-initial-gap}
Let \(I\in\mathcal I(\mathbf C,\Delta)\), and suppose that
\(\gamma L_{\calD,\bx}<\frac12\) and \(s>0\). Then, for every \(K\ge0\),
\begin{equation}
    \mathcal W_0-\mathcal W_K\le\Delta_{\mathcal W}
    :=\Delta+\frac{L_h^2}{\rho}+\frac{(G_f+G_{\calA}L_h)^2}{2\mu_\gamma}.
    \label{eq:b2-prime-initial-gap}
\end{equation}
\end{lemma}
}

\begin{theorem}[Iteration complexity of Fast LRALM]
\label{thm:b2-prime-complexity}
Let \(\mathbf C=(L_f,G_f,L_{\calA},G_{\calA},L_h)\in\RR^5_{\ge0}\) and
\(\Delta>0\). Fix \(\rho,\gamma,\tau>0\) such that
\[
    \gamma L_{\calD,\bx}\le\frac{1}{10},
    \qquad
    0<\tau\le\frac{\rho}{4},
    \qquad
    14\gamma\tau G_{\calA}^2\le1,
\]
\red{set
\(C_R:=\max\left\{8\sqrt{\gamma}\,L_{\calD,\bx}+\frac{6}{\sqrt{\gamma}},\,
\sqrt{\frac{8}{\tau}}\right\}\),}
and choose \(c_s>0\) with \(c_s\ge6C_\lambda L_h\). Given any
\(0<\epsilon\le c_s/\max\{\rho,4\tau\}\), set
\begin{equation}
    \eta=\gamma,\qquad s=\frac{c_s}{\epsilon},\qquad
    \alpha=\frac{1}{6\sqrt{\kappa}}=\frac{1}{6}\sqrt{\frac{\tau\epsilon}{c_s}}.
    \label{eq:b2-prime-accuracy-parameters}
\end{equation}
Then, for every \(I\in\mathcal I(\mathbf C,\Delta)\) and every integer
\[
    K\ge4C_R^2\Delta_{\mathcal W}\sqrt{\frac{c_s}{\tau}}\,
    \epsilon^{-5/2},
\]
at least one of the first \(K\) outputs of Algorithm~\ref{alg:b2-prime}
on \(I\) is an \(\epsilon\)-KKT point. With first-order information cached
between iterations, these \(K\) iterations and their outputs use \(3K+1\)
oracle calls; thus both the iteration and oracle complexities are
\(\mathcal O(\epsilon^{-5/2})\).
\end{theorem}
\begin{proof}
Since \(\epsilon\le c_s/\max\{\rho,4\tau\}\),
\eqref{eq:b2-prime-accuracy-parameters} gives \(s\ge\max\{\rho,4\tau\}\),
\(\kappa=s/\tau\ge4\), \(\tau L_\lambda\le\frac12\), and
\(\alpha=\frac{1}{6\sqrt{\kappa}}\le\frac{1}{12}\); hence~\eqref{eq:b2-prime-stepsizes}
holds. Let \(I\in\mathcal I(\mathbf C,\Delta)\), so
Assumptions~\ref{ass:smoothness} and~\ref{ass:h} hold with the constants in
\(\mathbf C\). Since \(\Delta_{\mathcal W}\ge\Delta>0\) and \(C_R>0\), the
lower bound on \(K\) implies \(K\ge1\).
Lemmas~\ref{lem:b2-prime-potential-descent} and~\ref{lem:b2-prime-initial-gap}
then yield
\[
    \sum_{k=0}^{K-1}(\mathcal W_k-\mathcal W_{k+1})
    =\mathcal W_0-\mathcal W_K\le\Delta_{\mathcal W}.
\]
Hence some \(0\le k<K\) satisfies
\(D_k:=\mathcal W_k-\mathcal W_{k+1}\le\Delta_{\mathcal W}/K\). At this
index \(k\), Lemma~\ref{lem:b2-prime-potential-descent} and
\(\alpha=\frac{1}{6\sqrt{\kappa}}\) give
\[
    \|\bx^{k+1}-\bx^k\|\le8\sqrt{2\gamma D_k},
    \qquad
    \|\blmbd^{k+1}-\bar{\blmbd}^k\|\le4\sqrt{\tau D_k},
    \qquad
    \|\bx^{k+1}-\bz^k\|\le6\sqrt{\gamma\sqrt{\kappa}\,D_k}.
\]
Substituting these bounds and Lemma~\ref{lem:b2-prime-multiplier-bounds} into
Lemma~\ref{lem:b2-prime-residual-decomposition} with \(\eta=\gamma\), and
using \(\sqrt{D_k}\le\sqrt{\sqrt{\kappa}\,D_k/2}\) from \(\kappa\ge4\), we
obtain
\[
    \min_{0\le k<K}\mathcal R_{\mathrm{KKT}}
    (\bx^{k+1},\by^{k+1},\bq^k)
    \le C_R\sqrt{\frac{\sqrt{\kappa}\,\Delta_{\mathcal W}}{K}}
    +\frac{3C_\lambda L_h}{s}.
\]
With \(\sqrt{\kappa}=\sqrt{c_s/(\tau\epsilon)}\), the lower bound on \(K\)
makes the first term at most \(\epsilon/2\), and
\(c_s\ge6C_\lambda L_h\) makes the second at most \(\epsilon/2\).
The constants \(C_R\), \(C_\lambda\), and \(\Delta_{\mathcal W}\)
are independent of \(\epsilon\).
Finally, with \(\mathcal O_{f,\calA}(\bx^k)\) cached between iterations,
each iteration uses one new call to \(\mathcal O_{f,\calA}\) and two calls
to \(\mathcal O_h\), in addition to the initial call
\(\mathcal O_{f,\calA}(\bx^0)\), giving \(3K+1\) oracle calls in total.
\end{proof}

